% !TEX root = main.tex
\section{Related Work}\label{related_work}

We briefly review related work concerning single-turn jailbreak attacks and multi-turn jailbreak attacks.

\textbf{Single-Turn Attacks:}
Early approaches~\cite{shen2023anything} relied on manually crafting prompts to execute jailbreak attacks. However, manual crafting was time and labor-intensive, leading attacks to gradually shift towards automation. The GCG method~\cite{zou2023universal} employed white-box attacks utilizing gradient information for jailbreak, yet resulting in poor readability of GCG-like outputs. AutoDAN~\cite{liu2023autodan} introduced genetic algorithms for automated updates, while Masterkey~\cite{deng2024masterkey} explored black-box approaches using time-based SQL injection to probe the defense mechanisms of LLM chatbots. Additionally, it leveraged fine-tuning and RFLH LLMs for automated jailbreak expansion. PAIR~\cite{chao2023jailbreaking} proposed iterative search in large model conversations, continuously optimizing single-turn attack prompts. GPTFUzz~\cite{yu2023gptfuzzer} combined attacks with fuzzing techniques, continually generating attack prompts based on template seeds. Furthermore, attacks such as multilingual attacks~\cite{deng2023multilingual} and obfuscation level attacks utilized low-resource training languages and instruction obfuscation\cite{shang2024can} to execute attacks.

However, single-turn jailbreak attack patterns are straightforward and thus easily detectable and defensible. As security alignments continue to strengthen, once the model is updated, previously effective prompts may become ineffective. Therefore, jailbreak attacks are now venturing towards multi-turn dialogues.

\textbf{Multi-Turn Jailbreak Attack:}
~\cite{li2023multi} employed multi-turn dialogues to carry out jailbreak attacks, circumventing the limitations of LLMs, presenting privacy and security risks, and extracting personally identifiable information (PII). ~\cite{zhou2024speak} utilized manual construction of multi-turn templates, harnessing GPT-4 for automated generation, to progressively intensify malicious intent and execute jailbreak attacks through sentence and goal reconstruction. ~\cite{russinovich2024great}facilitated benign interactions between large and target models, using the model’s own outputs to gradually steer the model in task execution, thereby achieving multi-turn jailbreak attacks. ~\cite{bhardwaj2023red}conducted an exploration of Conversation Understanding (CoU) prompt chains for jailbreak attacks on LLMs, alongside the creation of a red team dataset and the proposal of a security alignment method based on gradient ascent to penalize harmful responses. ~\cite{li2024drattack}decomposed original prompts into sub-prompts and subjected them to semantically similar but harmless implicit reconstruction, analyzing syntax to replace synonyms, thus preserving the original intent while undermining the security constraints of the language model. ~\cite{yang2024chain} proposed a semantic-driven context multi-turn attack method, adapting attack strategies adaptively through context feedback and semantic relevance in multi-turn dialogues of LLMs, thereby achieving semantic-level jailbreak attacks. Additionally, there are strategies that utilize multi-turn interactions to achieve puzzle games~\cite{DBLP:conf/uss/LiuZZDM024}, thus obscuring prompts and other non-semantic multi-turn jailbreak attack strategies.

Presently, multi-turn semantic jailbreak attacks exhibit vague strategies and high false positive rates. We attribute this to the unclear positioning of multi-turn interactions within jailbreaking and excessively complex strategies. Therefore, we have re-examined the advantages of multi-turn attacks and proposed a multi-turn contextual fusion attack strategy.


\textbf{Factors Influencing Jailbreak Attacks}
~\cite{Zou2024IsTS} delved into the impact of system prompts on prison prompts within LLM, revealing the transferable characteristics of prison prompts and proposing an evolutionary algorithm targeting system prompts to enhance the model’s robustness against them. ~\cite{qi2024finetuning} unveiled the security risks posed by fine-tuning LLMs, demonstrating that malicious fine-tuning can easily breach the model’s security alignment mechanism. ~\cite{huang2024catastrophic} discovered vulnerabilities in existing alignment procedures and assessments, which may be based on default decoding settings and exhibit flaws when configurations vary slightly.
~\cite{zhang2024large} demonstrated that even if LLM rejects toxic queries, harmful responses can be concealed within the top k hard label information, thereby coercing the model to divulge it during autoregressive output generation by enforcing the use of low-rank output tokens, enabling jailbreak attacks.

In some multi-turn approaches, merging multi-turn contexts is considered as one of the direct influencing factors in jailbreak attacks. However, in this paper, we argue that the context plays an indirectly supportive role rather than achieving the same impact level as system prompt templates.